% Copy this file with `make cover-init`, then edit only the ignored copy.
% Escape LaTeX special characters in prose, for example: R\&D, 50\%, and name\_id.

\newcommand{\applicantname}{Applicant Name}
\newcommand{\applicantlocation}{City, Country}
\newcommand{\applicantemail}{applicant@example.com}
\newcommand{\applicantwebsiteurl}{https://example.com}
\newcommand{\applicantwebsitelabel}{example.com}
\newcommand{\applicantgithuburl}{https://github.com/example}
\newcommand{\applicantgithublabel}{github.com/example}
\newcommand{\applicantlinkedinurl}{https://linkedin.com/in/example}
\newcommand{\applicantlinkedinlabel}{linkedin.com/in/example}

\newcommand{\applicationdate}{1 January 20XX}
\newcommand{\recipientline}{Hiring Team}
\newcommand{\salutation}{Hiring Team}
\newcommand{\companyname}{Example Company}
\newcommand{\positiontitle}{Example Position}

% Use complete paragraphs instead of small sentence fragments. This keeps the
% writing natural and makes every company-specific claim easy to review.
\newcommand{\openingparagraph}{Explain briefly why this role and company are a good match.}
\newcommand{\evidenceparagraph}{Give the strongest relevant evidence from recent work, using only claims that can be substantiated.}
\newcommand{\closingparagraph}{Connect that evidence to the role and state the contribution you would like to make.}
